% ============================================================
%  Inditex: liquidez, endeudamiento y flujo de caja (2022-2026)
%  Compilar con XeLaTeX (Overleaf: Menu > Compiler > XeLaTeX)
% ============================================================
\documentclass[11pt,a4paper]{article}

% ---------- Idioma y tipografía ----------
\usepackage{fontspec}            % Latin Modern por defecto
% \setmainfont{TeX Gyre Pagella} % alternativa: descomentar para Palatino
\usepackage{polyglossia}
\setdefaultlanguage{spanish}
\usepackage{etoolbox}
\gappto\captionsspanish{\renewcommand{\tablename}{Tabla}\renewcommand{\listtablename}{Índice de tablas}}

% ---------- Paquetes ----------
\usepackage[a4paper,top=2.8cm,bottom=2.6cm,left=2.7cm,right=2.7cm,headheight=14pt]{geometry}
\usepackage{xcolor}
\usepackage[most]{tcolorbox}
\usepackage{colortbl}
\usepackage{graphicx}
\usepackage{booktabs}
\usepackage{tabularx}
\usepackage{array}
\usepackage{multirow}
\usepackage{enumitem}
\usepackage{titlesec}
\usepackage{fancyhdr}
\usepackage{caption}
\usepackage{float}
\usepackage{microtype}
\usepackage[hidelinks]{hyperref}

% ---------- Formato de títulos ----------
% ---------- Colores ----------
\definecolor{navy}{HTML}{14284B}
\definecolor{accent}{HTML}{1F5FBF}
\definecolor{softgrey}{HTML}{F2F3F5}
\definecolor{rulegrey}{HTML}{C9CCD2}

\titleformat{\section}{\normalfont\Large\bfseries\color{navy}}{\thesection.}{0.7em}{}[\vspace{-0.35em}{\color{rulegrey}\titlerule[0.6pt]}]
\titleformat{\subsection}{\normalfont\large\bfseries\color{accent}}{\thesubsection.}{0.6em}{}
\titlespacing*{\section}{0pt}{2.0em}{0.9em}
\titlespacing*{\subsection}{0pt}{1.4em}{0.5em}

% ---------- Cabecera y pie ----------
\pagestyle{fancy}
\fancyhf{}
\fancyhead[L]{\small\color{navy}\textsc{Inditex · Análisis económico-financiero}}
\fancyhead[R]{\small\color{navy}\textsc{MEFI 2026--27}}
\fancyfoot[C]{\small\thepage}
\renewcommand{\headrulewidth}{0.4pt}
\renewcommand{\headrule}{{\color{rulegrey}\hrule height\headrulewidth width\headwidth}}
\fancypagestyle{plain}{\fancyhf{}\fancyfoot[C]{\small\thepage}\renewcommand{\headrulewidth}{0pt}}

% ---------- Párrafos, listas y leyendas ----------
\setlength{\parskip}{0.55em}
\setlength{\parindent}{0pt}
\setlength{\emergencystretch}{2.5em}
\linespread{1.08}
\setlist{itemsep=0.3em,topsep=0.4em,leftmargin=1.5em}
\captionsetup{font=small,labelfont={bf,color=navy},labelsep=period,margin=0.5cm}
\captionsetup[table]{position=top,skip=6pt}
\captionsetup[figure]{skip=6pt}

% ---------- Utilidades ----------
\newcommand{\mn}{\ensuremath{-}}                 % signo menos tipográfico
% Recuadro con la frase-resumen de cada apartado
\newtcolorbox{ideabox}{enhanced,colback=softgrey,colframe=softgrey,borderline west={3pt}{0pt}{accent},
  sharp corners,boxrule=0pt,left=10pt,right=10pt,top=7pt,bottom=7pt,fontupper=\itshape,before skip=0.6em,after skip=0.9em}
\newcommand{\idea}[1]{\begin{ideabox}#1\end{ideabox}}
\newcommand{\fuente}[1]{\par\vspace{2pt}{\footnotesize\textit{Fuente:} #1\par}}
\newcolumntype{R}{>{\raggedleft\arraybackslash}X}

% ============================================================
\begin{document}

% ---------- Portada ----------
\begin{titlepage}
  \centering
  \vspace*{1.2cm}
  {\large\color{navy}\textsc{Universidad Politécnica de Madrid}\par}
  \vspace{0.15cm}
  {\large\color{navy}\textsc{Instituto BME}\par}
  \vspace{0.6cm}
  {\normalsize Master in Financial Engineering and Innovation (MEFI)\par}
  {\normalsize Análisis de la Situación Económico-Financiera\par}

  \vspace{3.2cm}
  {\color{navy}\rule{\textwidth}{1.2pt}}\par
  \vspace{0.5cm}
  {\LARGE\bfseries\color{navy} Análisis financiero de Inditex\par}
  \vspace{0.35cm}
  {\Large\color{accent} Liquidez, endeudamiento y flujo de caja (2022--2026)\par}
  \vspace{0.5cm}
  {\color{navy}\rule{\textwidth}{0.4pt}}\par

  \vspace{3.2cm}
  {\large\color{accent}\textit{Autores}\par}
  \vspace{0.3cm}
  {\large Álvaro Pérez Ramos\par}
  \vspace{0.1cm}
  {\large Anxo Brandán Herrera Castro\par}
  \vspace{0.1cm}
  {\large Javier Lázaro Rodríguez\par}

  \vspace{1.4cm}
  {\normalsize\textit{Profesor:} Eduardo García-Mauriño\par}

  \vfill
  {\normalsize Madrid, octubre de 2026\par}
\end{titlepage}

% ---------- Índice ----------
\tableofcontents
\vspace{1.5em}
\listoffigures
\listoftables
\clearpage

% ============================================================
\section*{Resumen ejecutivo}
\addcontentsline{toc}{section}{Resumen ejecutivo}

\idea{Inditex se financia con su propio negocio: cobra antes de pagar, tiene más caja que deuda y
genera caja suficiente para invertir y pagar dividendos crecientes.}

Entre 2022 y 2026 las ventas crecen un 44\,\% (de 27.716 a 39.864 M€) y el EBITDA un 57\,\%
(de 7.183 a 11.267 M€) sin tensionar su posición financiera. La Tabla~\ref{tab:kpi} resume los
cuatro indicadores que mejor describen esa situación.

\begin{table}[H]
  \centering
  \caption{Indicadores clave del ejercicio 2026}
  \label{tab:kpi}
  \begin{tabular}{@{}l r l@{}}
    \toprule
    \color{navy}\textbf{Indicador} & \color{navy}\textbf{Valor} & \color{navy}\textbf{Lectura} \\
    \midrule
    Superávit de financiación (FM \mn{} NOF) & \textbf{\color{accent}9.376 M€} & Liquidez muy holgada \\
    Ciclo de caja (PMA + PMC \mn{} PMP)      & \textbf{\color{accent}\mn18,5 días} & Los proveedores financian el circulante \\
    Deuda neta                               & \textbf{\color{accent}\mn4.406 M€} & Caja neta, sin deuda bancaria \\
    Dividendos / cash-flow libre (2022--26)  & \textbf{\color{accent}74\,\%} & El dividendo absorbe casi todo el flujo libre \\
    \bottomrule
  \end{tabular}
\end{table}

\begin{enumerate}
  \item \textbf{Liquidez muy holgada.} Las NOF son negativas (\mn3.857 M€) y el Fondo de Maniobra positivo
        (5.519 M€): la empresa tiene un superávit de financiación a corto plazo de 9.376 M€.
  \item \textbf{Los proveedores financian el circulante.} Inditex cobra en unos 6 días, vende su stock en
        unos 71 y paga a unos 96: su ciclo de caja es de \mn18,5 días.
  \item \textbf{No tiene deuda neta, sino caja neta} (4.406 M€, 0,39 años de EBITDA). El EBITDA cubre
        30 veces los gastos financieros.
  \item \textbf{Genera mucha caja, pero el dividendo ya absorbe casi todo el flujo libre.} De 2022 a 2026
        el cash-flow libre acumulado es 25.646 M€ y los dividendos 18.882 M€. El dividendo se ha
        multiplicado por 2,4, más rápido que el beneficio, y en 2026 la caja neta baja 720 M€.
\end{enumerate}

{\small\textit{Nota:} los años corresponden a la fecha de cierre del ejercicio (31 de enero); «2026» es el
ejercicio fiscal de febrero de 2025 a enero de 2026. Todas las cifras se expresan en millones de euros (M€).}

% ============================================================
\section{Liquidez: ¿puede pagar sus obligaciones a corto plazo?}

\idea{Sí, y con mucho margen: el negocio de Inditex no necesita financiación para su circulante, sino que la genera.}

\subsection{Fondo de Maniobra frente a NOF}

En clase vimos dos ideas complementarias:
\begin{itemize}
  \item \textbf{Fondo de Maniobra (FM)} = capitales permanentes \mn{} activo no corriente. Es la
        financiación a largo plazo que \emph{sobra} después de pagar el inmovilizado y queda disponible
        para el circulante.
  \item \textbf{Necesidades Operativas de Fondos (NOF)} = existencias + clientes \mn{} proveedores (activo
        menos pasivo corriente operativo). Es lo que la operativa diaria \emph{necesita} que se financie.
\end{itemize}

Si FM > NOF hay \textbf{superávit} y el exceso se acumula como tesorería. Si FM < NOF hay
\textbf{déficit} y hay que pedir crédito a corto plazo.

En Inditex las NOF son \textbf{negativas todos los años} (\mn3.857 M€ en 2026, unos 35 días de ventas).
Su operativa no consume fondos, los aporta: cuando vende una prenda, todavía no ha pagado al proveedor
que la fabricó. Como además el FM es positivo, el superávit (FM \mn{} NOF) ronda los 9.400 M€.
Coincide exactamente con su tesorería neta (caja + inversiones financieras a corto \mn{} deuda a corto),
que es la comprobación de la ecuación del equilibrio financiero (Figura~\ref{fig:fm}).

\begin{figure}[H]
  \centering
  \includegraphics[width=\textwidth]{figuras/fig1_fm_nof.pdf}
  \caption[Fondo de Maniobra, NOF y superávit de financiación (2022--2026)]{Fondo de Maniobra, NOF y superávit de financiación (M€, a 31 de enero). Aunque el FM cae desde 2024, el superávit se mantiene por encima de 9.300 M€.}
  \label{fig:fm}
\end{figure}

Desde 2024 el FM cae de 7.079 a 5.519 M€. La razón es que Inditex invierte más en activo fijo
(el capex pasa de 1.126 a 2.712 M€, sobre todo logística y tiendas) y reparte más dividendos.
Aun así, el superávit apenas se mueve, porque al mismo tiempo las NOF se hacen más negativas:
lo que pierde por el FM lo recupera con la financiación de proveedores.

\subsection{El ciclo de caja: cobra antes de pagar}

Los periodos medios explican por qué las NOF son negativas. En 2026, Inditex tarda unos
\textbf{71 días} en vender su stock (PMA) y unos \textbf{6 días} en cobrar (PMC); casi todas sus ventas se
cobran al contado en tienda u online, y el saldo de clientes corresponde sobre todo a franquicias.
En cambio, paga a sus proveedores a unos \textbf{96 días} (PMP).

\begin{figure}[H]
  \centering
  \includegraphics[width=\textwidth]{figuras/fig2_ciclo_caja.pdf}
  \caption[Ciclo de caja del ejercicio 2026]{Ciclo de caja del ejercicio 2026, en días. Inditex cobra la venta unas dos semanas y media antes de pagar al proveedor.}
  \label{fig:ciclo}
\end{figure}

El ciclo de caja (PMA + PMC \mn{} PMP) es de \textbf{\mn18,5 días}: el dinero de la venta entra unas
dos semanas y media antes de tener que pagar la mercancía (Figura~\ref{fig:ciclo}). El PMA mejora de 93 a
71 días, lo que indica una mejor rotación del inventario. El PMP baja de 111 a 96 días: los proveedores
financian algo menos, y es un punto a vigilar.

\subsection{Ratios de liquidez}

La liquidez general (AC/PC) es \textbf{1,54}, por debajo de la «regla del 2», pero en este caso no es
una señal de alarma. Como vimos en clase, el ratio depende del modelo de negocio: el pasivo corriente
de Inditex es sobre todo deuda con proveedores, que no tiene coste y se renueva sola. Lo relevante es
que el test ácido (1,22) y las disponibilidades líquidas (1,08) están por encima de 1. Solo con su caja
e inversiones financieras a corto, Inditex podría pagar todo su pasivo corriente.

% ============================================================
\section{Endeudamiento: ¿está Inditex endeudada?}

\idea{No. Inditex tiene 4.406 M€ más en activos líquidos que en deuda financiera, es decir, caja neta.}

La deuda neta es la deuda financiera bruta menos los activos tesoreros (caja e inversiones
financieras a corto). En 2026 la deuda bruta es de 6.554 M€ y los activos tesoreros suman
10.960 M€, por lo que la deuda neta sale \textbf{negativa}: \mn4.406 M€. Además, cerca del 90\,\% de esa
deuda bruta son arrendamientos de tiendas registrados como deuda por la NIIF 16; la deuda con bancos es
de apenas 2 M€.

Los ratios que vimos en clase lo confirman:
\begin{itemize}
  \item \textbf{Deuda neta / EBITDA = \mn0,39.} Los bancos empiezan a preocuparse por encima de 3 veces.
        En Inditex el ratio es negativo: con la caja actual podría devolver toda su deuda y aún le sobraría
        el equivalente a unos 5 meses de EBITDA.
  \item \textbf{EBITDA / gastos financieros = 30 veces.} Baja desde 64 en 2022 porque suben los intereses de
        los arrendamientos y los tipos, pero sigue siendo una cobertura enorme.
  \item \textbf{Solvencia (activo / pasivo) = 2,34} y \textbf{pasivo / patrimonio neto = 0,75.} Los bienes de la
        empresa cubren más de dos veces sus deudas, y los fondos propios superan al pasivo total.
\end{itemize}

La caja neta alcanza su máximo en 2024 (5.220 M€) y desde entonces baja unos 800 M€. No es por pedir
prestado, sino porque los dividendos y el capex crecen más rápido que la caja generada, como se ve
en el apartado siguiente.

% ============================================================
\section{Flujo de caja: ¿de dónde sale el dinero y a dónde va?}

\idea{Entre 2022 y 2026, Inditex generó 40.614 M€ de caja con su operativa. Después de invertir le
quedaron 25.646 M€ de flujo libre, que destinó sobre todo a dividendos (18.882 M€).}

El Estado de Flujos de Efectivo ordena los movimientos en una cascada, como vimos en el tema 4:
recursos generados por las operaciones, variación de las NOF, cash-flow operativo, inversiones,
cash-flow libre, pagos al accionista, variación de la deuda y, por último, variación de caja.
La Figura~\ref{fig:cascada} muestra esa cascada con los cinco años acumulados.

\begin{figure}[H]
  \centering
  \includegraphics[width=\textwidth]{figuras/fig3_cascada.pdf}
  \caption[Cascada del flujo de caja, acumulado 2022--2026]{Del recurso generado a la variación de caja, acumulado 2022--2026 (M€). El dividendo se lleva tres de cada cuatro euros de cash-flow libre.}
  \label{fig:cascada}
\end{figure}

Tres lecturas clave:
\begin{itemize}
  \item \textbf{Beneficio de calidad.} El cash-flow operativo de 2026 (9.232 M€) supera al beneficio neto
        (6.220 M€): el beneficio se convierte en caja. La diferencia se debe sobre todo a las amortizaciones,
        que son gasto pero no salida de dinero.
  \item \textbf{No todo el flujo de inversión es inversión productiva.} El capex real (tiendas, logística,
        tecnología) se duplica, de 1.126 a 2.712 M€. Pero de los 14.968 M€ de inversión acumulada, 5.489 M€
        son excedentes colocados en inversiones financieras temporales, que no es inversión sino gestión de
        tesorería: el dinero sigue siendo de Inditex.
  \item \textbf{La «variación de deuda» son, en realidad, alquileres.} Los \mn8.664 M€ que el EFE clasifica
        como devolución de deuda son los pagos por arrendamientos de tiendas (NIIF 16), un coste recurrente
        del negocio.
\end{itemize}

Si se suma todo, el colchón se estrecha. Al flujo libre (sin las colocaciones de tesorería) se le restan
los alquileres y los dividendos. El resultado pasa de +1.840 M€ en 2022 a \mn508 M€ en 2026: hoy el
dividendo consume todo lo que queda después de invertir y pagar los alquileres. Por eso la caja
neta baja 720 M€ en 2026.

% ============================================================
\section{Conclusiones}

\idea{Inditex es un ejemplo de manual de empresa que crece sin necesitar financiación externa: su modelo de
negocio genera la liquidez que consume.}

\begin{itemize}
  \item \textbf{Fortaleza estructural.} Tiene NOF negativas, un ciclo de caja negativo y caja neta. Cuanto más
        vende, más financiación gratuita obtiene de sus proveedores. Es el caso contrario al de una empresa
        industrial típica, que necesita crédito para financiar su circulante.
  \item \textbf{Riesgo financiero mínimo.} No tiene deuda bancaria y sus coberturas son muy holgadas. Los ratios
        «malos» en apariencia, como una liquidez inferior a 2 o un pasivo que crece, se explican por la
        financiación de proveedores y los alquileres, no por deuda con coste.
  \item \textbf{Punto de atención: el destino de la caja.} El capex y, sobre todo, el dividendo crecen más rápido
        que la caja generada. Por eso el FM y la caja neta bajan desde 2024. Hoy no es un problema, porque
        sigue habiendo 4.406 M€ de caja neta. Pero si esa política se mantiene, Inditex iría consumiendo su
        colchón de liquidez.
\end{itemize}

En definitiva, Inditex presenta una \textbf{liquidez y una solvencia excelentes}, con una política de
retribución al accionista que ya reparte prácticamente todo lo que el negocio genera.

% ============================================================
\clearpage
\appendix
\renewcommand{\thesection}{\Alph{section}}
\section{Anexo: ratios calculados}

Se recogen a continuación todos los ratios calculados en el libro de trabajo, se comenten o no en el
cuerpo del informe. Cifras en M€ salvo indicación; los años corresponden a la fecha de cierre del
ejercicio (31 de enero).

\begin{table}[H]
\centering
\caption{Ratios de liquidez}
\label{tab:liq}
\small
\begin{tabularx}{\textwidth}{@{}>{\raggedright\arraybackslash}p{6.3cm}RRRRR@{}}
\toprule
\color{navy}\textbf{Concepto} & \color{navy}\textbf{2022} & \color{navy}\textbf{2023} & \color{navy}\textbf{2024} & \color{navy}\textbf{2025} & \color{navy}\textbf{2026} \\
\midrule
Fondo de Maniobra & 5.573 & 6.502 & 7.079 & 6.169 & 5.519 \\
NOF & \mn2.203 & \mn2.005 & \mn2.873 & \mn3.736 & \mn3.857 \\
Superávit (+) / déficit (\mn) = FM \mn{} NOF & 7.776 & 8.507 & 9.952 & 9.905 & 9.376 \\
Tesorería neta & 7.776 & 8.507 & 9.952 & 9.905 & 9.376 \\
Variación del FM & -- & 929 & 577 & \mn910 & \mn650 \\
Variación de las NOF & -- & 198 & \mn868 & \mn863 & \mn121 \\
Variación de la liquidez & -- & 731 & 1.445 & \mn47 & \mn529 \\
\midrule
Liquidez general (AC/PC) & 1,69 & 1,80 & 1,79 & 1,61 & 1,54 \\
Test ácido ((AC \mn{} existencias)/PC) & 1,32 & 1,41 & 1,46 & 1,28 & 1,22 \\
Disponibilidades líquidas ((caja + IF c/p)/PC) & 1,17 & 1,24 & 1,28 & 1,13 & 1,08 \\
\midrule
PMC (días) & 5,5 & 5,5 & 5,2 & 5,8 & 6,0 \\
PMA (días) & 93,3 & 83,1 & 71,3 & 74,4 & 71,3 \\
PMP (días) & 110,8 & 96,8 & 102,6 & 99,9 & 95,7 \\
Ciclo de caja PMA + PMC \mn{} PMP (días) & \mn12,0 & \mn8,2 & \mn26,1 & \mn19,7 & \mn18,5 \\
\bottomrule
\end{tabularx}
\end{table}

\begin{table}[H]
\centering
\caption{Ratios de endeudamiento}
\label{tab:end}
\small
\begin{tabularx}{\textwidth}{@{}>{\raggedright\arraybackslash}p{6.3cm}RRRRR@{}}
\toprule
\color{navy}\textbf{Concepto} & \color{navy}\textbf{2022} & \color{navy}\textbf{2023} & \color{navy}\textbf{2024} & \color{navy}\textbf{2025} & \color{navy}\textbf{2026} \\
\midrule
Deuda financiera bruta & 6.417 & 6.004 & 6.202 & 6.376 & 6.554 \\
Activos tesoreros & 9.395 & 10.083 & 11.422 & 11.502 & 10.960 \\
Deuda neta (negativo = caja neta) & \mn2.978 & \mn4.079 & \mn5.220 & \mn5.126 & \mn4.406 \\
Variación de la deuda neta & -- & \mn1.101 & \mn1.141 & 94 & 720 \\
\midrule
EBITDA / gastos financieros (x) & 63,6 & 60,1 & 36,3 & 30,4 & 30,5 \\
Deuda neta / EBITDA (x) & \mn0,41 & \mn0,47 & \mn0,53 & \mn0,48 & \mn0,39 \\
Deuda neta / flujo de caja orgánico (x) & \mn0,52 & \mn0,77 & \mn0,76 & \mn0,76 & \mn0,67 \\
Apalancamiento neto DN/(DN + PN) & \mn0,23 & \mn0,31 & \mn0,39 & \mn0,35 & \mn0,28 \\
Solvencia (activo / pasivo) & 2,20 & 2,32 & 2,33 & 2,31 & 2,34 \\
Endeudamiento bruto (pasivo / PN) & 0,84 & 0,76 & 0,75 & 0,76 & 0,75 \\
Endeudamiento neto (DN / PN) & \mn0,19 & \mn0,24 & \mn0,28 & \mn0,26 & \mn0,22 \\
\bottomrule
\end{tabularx}
\end{table}

\begin{table}[H]
\centering
\caption{Análisis del flujo de caja (M€)}
\label{tab:fc}
\small
\begin{tabularx}{\textwidth}{@{}>{\raggedright\arraybackslash}p{4.7cm}RRRRRR@{}}
\toprule
\color{navy}\textbf{Concepto} & \color{navy}\textbf{2022} & \color{navy}\textbf{2023} & \color{navy}\textbf{2024} & \color{navy}\textbf{2025} & \color{navy}\textbf{2026} & \color{navy}\textbf{2022--26}\hspace{-2pt} \\
\midrule
a) Recursos generados por las operaciones & 6.532 & 7.342 & 8.724 & 9.486 & 10.034 & 42.119 \\
(3) Variación NOF & 223 & \mn669 & \mn57 & \mn199 & \mn803 & \mn1.505 \\
\textbf{I) Cash-flow operativo} & \textbf{6.755} & \textbf{6.673} & \textbf{8.667} & \textbf{9.288} & \textbf{9.232} & \textbf{40.614} \\
II) Cash-flow de inversiones & \mn3.253 & \mn3.503 & \mn1.709 & \mn3.288 & \mn3.215 & \mn14.968 \\
\textbf{III) Cash-flow libre} & \textbf{3.502} & \textbf{3.170} & \textbf{6.958} & \textbf{6.000} & \textbf{6.016} & \textbf{25.646} \\
(9) Variación de patrimonio & \mn71 & \mn61 & 0 & 0 & 0 & \mn132 \\
(11) Dividendos & \mn2.192 & \mn2.914 & \mn3.744 & \mn4.797 & \mn5.235 & \mn18.882 \\
\textbf{IV) Cash-flow generado internamente} & \textbf{1.239} & \textbf{195} & \textbf{3.214} & \textbf{1.203} & \textbf{782} & \textbf{6.632} \\
(10) Variación de deuda (arrendamientos) & \mn1.646 & \mn1.638 & \mn1.730 & \mn1.810 & \mn1.839 & \mn8.664 \\
(D) Efecto tipo de cambio & 31 & \mn17 & \mn38 & \mn18 & \mn48 & \mn90 \\
\textbf{V) Variación de caja} & \textbf{\mn377} & \textbf{\mn1.460} & \textbf{1.446} & \textbf{\mn625} & \textbf{\mn1.106} & \textbf{\mn2.122} \\
\midrule
Tesorería final & 7.021 & 5.561 & 7.007 & 6.382 & 5.276 & -- \\
\bottomrule
\end{tabularx}
\end{table}

\subsection*{Notas metodológicas}
\begin{itemize}
  \item \textbf{NOF} = existencias + deudores comerciales + activos por impuesto corriente + otros activos
        corrientes + otros activos financieros + activos mantenidos para la venta \mn{} acreedores comerciales.
  \item \textbf{Periodos medios:} PMC = (deudores comerciales + franquicias) / (ventas × 1,21 / 365);
        PMP = proveedores / (compras × 1,21 / 365); PMA = existencias / (coste de ventas / 365).
        Se usan saldos de cierre y un IVA del 21\,\% como simplificación.
  \item \textbf{Deuda neta} = deuda financiera bruta (entidades de crédito, arrendamientos, provisiones y otros
        pasivos) \mn{} (efectivo + inversiones financieras a corto). Un valor negativo indica caja neta.
  \item \textbf{Flujo de caja orgánico} = cash-flow libre sin las colocaciones en inversiones financieras temporales.
  \item \textbf{Fuente:} cuentas anuales consolidadas de Inditex (NIIF-UE), con los datos del libro de trabajo
        facilitado en la asignatura.
\end{itemize}

\end{document}
